% Figure — Contre-réaction tension-parallèle : amplificateur A et réseau de contre-réaction β_F (EANA3, chapitre 2).
% D'après les notes externes [notes-2023-c04-06, p. 6], redessinée.
% Rendu SVG :
%   latex -output-directory=<dossier temporaire> 02-cr-tension-parallele.tex
%   dvisvgm --no-fonts --exact-bbox -o 02-cr-tension-parallele.svg <dossier temporaire>/02-cr-tension-parallele.dvi
\documentclass[border=4pt]{standalone}
\usepackage[T1]{fontenc}
\usepackage[utf8]{inputenc}
\usepackage{amsmath}
\usepackage{sansmath}
\renewcommand{\familydefault}{\sfdefault}
\usepackage[european, straightvoltages, americancurrents]{circuitikz}
\begin{document}
\sansmath
\begin{circuitikz}[font=\small, >=latex, fleche/.style={->, thick}]
  \draw[thick, fill=black!3] (1.9,1.6) rectangle (7.4,4.3); \node[font=\bfseries] at (7.1,4.05) {$A$};
  \draw[thick, fill=black!3] (1.9,-1.7) rectangle (7.4,0.6); \node at (4.8,-1.45) {$\beta_F$};
  \draw (1.9,3.5) node[circ]{} -- (2.8,3.5) to[generic, l=$Z_e$] (2.8,2.0) -- (1.9,2.0) node[circ]{};
  \draw[fleche] (2.25,2.15) -- node[left]{$v_e$} (2.25,3.35);
  \draw (4.6,3.5) to[american controlled voltage source, l=$A\,i_e$] (4.6,2.0);
  \draw (4.6,3.5) to[generic, l=$Z_s$] (6.6,3.5) -- (7.4,3.5) node[circ]{};  \draw (4.6,2.0) -- (7.4,2.0) node[circ]{};
  \draw (1.9,0.2) node[circ]{} -- (3.0,0.2) to[american controlled current source, l=$\beta_F\, v_o$] (3.0,-1.2) -- (1.9,-1.2) node[circ]{};
  \draw (7.4,0.2) node[circ]{}; \draw (7.4,-1.2) node[circ]{};
  \draw[fleche] (7.0,-1.05) -- node[left]{$v_o$} (7.0,0.05);
  \draw (0,3.5) node[ocirc]{}; \draw (0,-1.2) node[ocirc]{}; \draw (9.2,3.5) node[ocirc]{}; \draw (9.2,-1.2) node[ocirc]{};
  \draw[fleche] (-0.35,-1.0) -- node[left]{$v_i$} (-0.35,3.3);
  \draw[fleche] (9.55,-1.0) -- node[right]{$v_o$} (9.55,3.3);
  \draw (0,3.5) -- (1.9,3.5);  \draw (0.8,3.5) node[circ]{} -- (0.8,0.2) -- (1.28,0.2) arc[start angle=180, end angle=0, radius=0.12] -- (1.9,0.2);
  \draw (1.9,2.0) -- (1.4,2.0) -- (1.4,-1.2) node[circ]{};  \draw (0,-1.2) -- (1.9,-1.2);
  \draw[fleche] (0.1,3.7) -- node[above]{$i_i$} (0.7,3.7);  \draw[fleche] (1.0,3.7) -- node[above]{$i_e$} (1.7,3.7);
  \draw[fleche] (0.6,2.4) -- node[left]{$i_{fb}$} (0.6,1.2);
  \draw (7.4,3.5) -- (9.2,3.5);  \draw (8.5,3.5) node[circ]{} -- (8.5,0.2) -- (8.02,0.2) arc[start angle=0, end angle=180, radius=0.12] -- (7.4,0.2);
  \draw (7.4,2.0) -- (7.9,2.0) -- (7.9,-1.2) node[circ]{};  \draw (7.4,-1.2) -- (9.2,-1.2);
  \draw[fleche] (9.1,3.7) -- node[above]{$i_o$} (8.7,3.7);
\end{circuitikz}
\end{document}
